\documentclass[12pt]{article}
\input{structure.tex}
\renewcommand{\bottomtitlespace}{3cm}

\begin{document}

\renewcommand{\tl}{$2.2$ сек.}
\renewcommand{\ml}{$512$ MB}
\problem{Задача A3. NMST}
\addauthor{Автори: Христо Венев, Емил Инджев}{2025}{06}{07}
\renewcommand{\header}{
    НАЦИОНАЛЕН ЛЕТЕН ТУРНИР ПО ИНФОРМАТИКА\\
    Пловдив, 6-8 юни 2025 г.\\
    Група A, 11-12 клас
}

Даден е претеглен свързан граф с $N$ върха и $M$ ребра.
Графът е специален с това, че за всяко възможно тегло, най-много $R$ ребра имат това тегло.
Трябва да изчислите колко на брой минимални покриващи дървета има графът.

Покриващо дърво на граф е подмножество от ребрата му, образуващо дърво, включващо всички върхове.
Минимално покриващо дърво е покриващо дърво с минимална сума на теглата на избраните ребра.
Две покриващи дървета се считат за различни, ако множествата от ребра, които ги образуват, са различни.

\subsection{Вход}
На първия ред на стандартния вход се въвеждат естествените числа $N$ и $M$, отговарящи съответно на броя върхове и ребра.
На следващите $M$ реда се въвеждат тройки естествени числа $A_i$, $B_i$ и $C_i$ описващи ребрата -- ребро $i$ свързва върховете $A_i$ и $B_i$ и има тегло $C_i$.

\subsection{Изход}
Изведете броят покриващи дървета по модул $10^9+7$.

\subsection{Ограничения}
\begin{itemize}
\item $2 \leq N \leq 10^5$
\item $1 \leq M \leq 2 \times 10^5$
\item $1 \leq R \leq 16$
\item $1 \leq A_i < B_i \leq N$
\item $0 \leq C_i \leq 10^6$
\item $A_i \neq A_j$ или $B_i \neq B_j$, когато $i \neq j$
\end{itemize}

\subsection{Подзадачи}
\begin{table}[H]
    \begin{tblr}{|Q[c,m]|Q[c,m]|Q[c,m]|Q[c,m]|Q[c,m]|}
        \hline
        \textbf{Подзадача} & \textbf{Точки} & $N$ & $M$ & $R$ \\
        \hline
        $1$ &  $1$ & $\leq 10^5$ & $\leq 2 \times 10^5$ & $\leq 1$ \\
        \hline
        $2$ & $13$ & $\leq 10$ & $\leq 20$ & $\leq 10$ \\
        \hline
        $3$ & $27$ & $\leq 2000$ & $\leq 4000$ & $\leq 10$ \\
        \hline
        $4$ & $20$ & $\leq 10^5$ & $\leq 2 \times 10^5$ & $\leq 10$ \\
        \hline
        $5$ & $17$ & $\leq 10^5$ & $\leq 2 \times 10^5$ & $\leq 12$ \\
        \hline
        $6$ & $17$ & $\leq 10^5$ & $\leq 2 \times 10^5$ & $\leq 14$ \\
        \hline
        $7$ &  $5$ & $\leq 10^5$ & $\leq 2 \times 10^5$ & $\leq 16$ \\
        \hline
    \end{tblr}
    \caption*{Точките за дадена подзадача се получават само ако се преминат успешно всички тестове в нея и всички други подзадачи съдържащи се в нея.}
\end{table}

\newpage

\subsection{Пример}
\begin{table}[H]
    \begin{tblr}{|l|l|X[j]|}
        \hline
        \textbf{Вход} & \textbf{Изход} & \textbf{Обяснение на примера} \\
        \hline
        \texttt{\makecell[lt]{4 5 \\
        1 2 1 \\
        2 3 1 \\
        1 3 1 \\
        1 4 2 \\
        2 4 3}}
        &
        \texttt{3}
        &
        \makecell[lt]{Има три минимални покриващи дървета (с тегло 4): \\
        $\left\{(1,2), (2,3), (1,4)\right\}$ \\
        $\left\{(1,2), (1,3), (1,4)\right\}$ \\
        $\left\{(2,3), (1,3), (1,4)\right\}$ \\
        }\\
        \hline
        \texttt{\makecell[lt]{6 9 \\
        1 2 2 \\
        2 3 2 \\
        3 4 2 \\
        4 5 2 \\
        5 6 2 \\
        1 6 2 \\
        2 4 1 \\
        4 6 1 \\
        2 6 1}}
        &
        \texttt{24}
        &
        {Има 24 минимални покриващи дървета (с тегло 8).} \\
        \hline
    \end{tblr}
\end{table}

\end{document}
